\hypertarget{chapter-1-introduction}{%
\chapter{Chapter 1: Introduction}\label{chapter-1-introduction}}

\hypertarget{the-integration-challenge}{%
\section{1.1 The Integration
Challenge}\label{the-integration-challenge}}

\textbf{Systems biology aims to understand living organisms as
integrated systems}, not collections of isolated parts. A bacterial cell
responding to glucose depletion, for instance, involves: -
\textbf{Metabolic changes}: Enzyme kinetics shift as substrate
concentrations change (milliseconds to seconds) - \textbf{Gene
expression responses}: Transcription factors activate alternative
pathway genes (minutes to hours) - \textbf{Regulatory feedback}:
Metabolites regulate their own synthesis genes (closing the loop)

\textbf{Current computational models address these layers separately}: -
\textbf{Metabolic models} (e.g., flux balance analysis, SBML-encoded
pathways) capture enzyme kinetics but omit gene regulation -
\textbf{Gene regulatory networks} (Boolean logic, differential
equations) model transcription factors but ignore metabolite levels -
\textbf{Signaling pathway models} represent protein interactions but
lack biochemical context

\textbf{The result is a fragmented view}: Researchers build separate
models for metabolism, regulation, and signaling, then struggle to
integrate them. As Kitano (2002) argued, ``to understand biology at the
system level, we must examine the structure and dynamics of cellular and
organismal function, rather than the characteristics of isolated
parts.''

\textbf{The gap}: \textbf{No unified formalism spans metabolic
biochemistry and gene regulatory logic in a single compositional
framework.} Existing approaches either: 1. Model metabolism without
regulation (missing biological control) 2. Model regulation abstractly
without biochemical grounding (no mass balance) 3. Combine formalisms
ad-hoc (software integration, not formal semantics)

This thesis addresses the gap by \textbf{extending Petri net formalism}
to enable integrated multi-scale biological modeling.

\begin{center}\rule{0.5\linewidth}{0.5pt}\end{center}

\hypertarget{motivating-example-the-lac-operon-system}{%
\section{1.2 Motivating Example: The Lac Operon
System}\label{motivating-example-the-lac-operon-system}}

\hypertarget{biological-background}{%
\subsection{1.2.1 Biological Background}\label{biological-background}}

The \textbf{lactose operon} (\emph{lac} operon) in \emph{Escherichia
coli} is a classic example of gene regulation coupled to metabolism.
When lactose is present but glucose is absent, the cell expresses
enzymes to metabolize lactose. \textbf{This behavior requires
coordinating three layers}:

\textbf{1. Metabolic layer} (biochemical reactions): - \textbf{Glucose
metabolism}: Glucose \textrightarrow Pyruvate (glycolysis, 10 enzymatic steps) -
\textbf{Lactose metabolism}: Lactose + H₂O \textrightarrow Glucose + Galactose
(β-galactosidase enzyme) - \textbf{cAMP production}: When glucose is
low, adenylate cyclase produces cAMP from ATP

\textbf{2. Regulatory layer} (transcription factors): - \textbf{CRP
protein} (catabolite repressor protein): Binds cAMP \textrightarrow cAMP-CRP complex
(activator) - \textbf{LacI repressor}: Blocks \emph{lac} promoter unless
allolactose (lactose metabolite) binds - \textbf{Lac promoter}: Requires
\textbf{both} cAMP-CRP binding \textbf{and} LacI derepression for
transcription

\textbf{3. Integration} (metabolite-gene coupling): - \textbf{Metabolite
\textrightarrow Transcription}: Low glucose \textrightarrow High cAMP \textrightarrow cAMP-CRP \textrightarrow \emph{lacZ}
transcription - \textbf{Product feedback}: Lactose \textrightarrow Allolactose \textrightarrow LacI
derepression \textrightarrow More β-galactosidase

\hypertarget{why-this-cannot-be-modeled-with-existing-formalisms}{%
\subsection{1.2.2 Why This Cannot Be Modeled with Existing
Formalisms}\label{why-this-cannot-be-modeled-with-existing-formalisms}}

\textbf{Classical Petri nets} model biochemical reactions (glucose \textrightarrow
pyruvate) but lack: - \textbf{Regulatory arcs}: Cannot represent
``cAMP-CRP activates transcription'' as an arc (would require code) -
\textbf{Threshold logic}: Cannot encode ``LacI blocks promoter when
M(allolactose) \textless{} 0.1 mM'' in topology - \textbf{Heterogeneous
dynamics}: Glycolysis is continuous (ODE), gene expression is stochastic
bursts (Gillespie) \textrightarrow cannot mix in one model

\textbf{SBML} (Systems Biology Markup Language) models metabolic
reactions but: - \textbf{Gene regulation requires separate models}:
SBML-qual (qualitative) extension for Boolean logic, disconnected from
kinetics - \textbf{No compositional integration}: Metabolic SBML +
regulatory SBML-qual = two files, manually synchronized

\textbf{Boolean gene regulatory networks} model \emph{lacI} \textrightarrow
\emph{lacZ} logic but: - \textbf{No metabolite states}: Cannot represent
cAMP concentration (continuous variable) - \textbf{No stoichiometry}:
Cannot track lactose consumption, glucose production - \textbf{No enzyme
kinetics}: β-galactosidase activity is binary (on/off), not
Michaelis-Menten

\textbf{ODE systems} can model both layers by writing coupled
differential equations:

\begin{lstlisting}
d[Glucose]/dt = -v_glycolysis
d[cAMP]/dt = k_synth · (1 - [Glucose]/K) - k_deg · [cAMP]
d[lacZ_mRNA]/dt = k_tx · [cAMP-CRP] · (1 - [LacI]) - k_deg · [lacZ_mRNA]
\end{lstlisting}

\textbf{But}: - \textbf{Not compositional}: Adding a new metabolite or
gene requires rewriting equations (no modularity) - \textbf{No visual
topology}: Regulatory connections hidden in equation code - \textbf{No
formal analysis}: Cannot use Petri net structural analysis
(P-invariants, reachability)

\hypertarget{what-is-needed}{%
\subsection{1.2.3 What Is Needed}\label{what-is-needed}}

\textbf{A unified formalism must support}: 1. \textbf{Metabolic
reactions} with stoichiometry (Glucose + ATP \textrightarrow G6P + ADP) 2.
\textbf{Gene expression} with regulatory logic (\emph{lacZ}
transcription requires cAMP-CRP and ¬LacI) 3. \textbf{Metabolite-gene
coupling} (cAMP concentration \textrightarrow transcription factor activity) 4.
\textbf{Visual topology} (arcs encode regulation, not hidden in code) 5.
\textbf{Heterogeneous dynamics} (continuous glycolysis + stochastic
transcription bursts) 6. \textbf{Mass balance} (atoms conserved: C₆H₁₂O₆
\textrightarrow 2 C₃H₄O₃)

\textbf{This thesis presents such a formalism}: Extended Biological
Petri Nets.

\begin{center}\rule{0.5\linewidth}{0.5pt}\end{center}

\hypertarget{research-questions}{%
\section{1.3 Research Questions}\label{research-questions}}

This thesis addresses six fundamental questions about integrated
biological modeling:

\hypertarget{rq1-cooperative-parallelism-weak-independence}{%
\subsection{RQ1: Cooperative Parallelism (Weak
Independence)}\label{rq1-cooperative-parallelism-weak-independence}}

\textbf{Can Petri nets support parallel execution when transitions share
places?}

\begin{itemize}
\tightlist
\item
  \textbf{Classical Petri net theory}: Transitions are \textbf{strongly
  independent} if they share no places (disjoint neighborhoods) \textrightarrow can
  fire simultaneously
\item
  \textbf{Biological reality}: Enzymes catalyze multiple reactions
  (shared catalyst place), pathways converge (shared product place)
\item
  \textbf{Example}: Hexokinase catalyzes glucose phosphorylation AND
  fructose phosphorylation \textrightarrow shares enzyme place
\item
  \textbf{Question}: Can we define \textbf{weak independence} allowing
  shared catalysts and products while preserving correctness?
\end{itemize}

\textbf{Hypothesis}: Transitions with \textbf{disjoint input places} (no
resource conflicts) but shared outputs or catalysts are \textbf{weakly
independent} \textrightarrow can execute in parallel without violating reachability.

\hypertarget{rq2-heterogeneous-transition-types}{%
\subsection{RQ2: Heterogeneous Transition
Types}\label{rq2-heterogeneous-transition-types}}

\textbf{Can continuous, stochastic, timed, and burst transitions coexist
in a single model with consistent semantics?}

\begin{itemize}
\tightlist
\item
  \textbf{Biological systems exhibit multiple temporal scales}:

  \begin{itemize}
  \tightlist
  \item
    \textbf{Continuous}: Enzyme kinetics (Michaelis-Menten), equilibria
    (ms-s)
  \item
    \textbf{Stochastic}: Gene expression bursts, rare molecular events
    (low copy numbers)
  \item
    \textbf{Timed}: Cell cycle checkpoints, scheduled events
  \item
    \textbf{Burst}: Transcriptional pulsing (chromatin remodeling)
  \end{itemize}
\item
  \textbf{Example}: Energy sensing motif requires continuous glycolysis
  + stochastic gene transcription bursts
\item
  \textbf{Question}: Can these four transition types fire in a unified
  hybrid simulation without semantic conflicts?
\end{itemize}

\textbf{Hypothesis}: A \textbf{hybrid scheduler} coordinating four
specialized engines (ODE, Gillespie, event queue, burst sampler) can
simulate heterogeneous transitions with provable correctness.

\hypertarget{rq3-arc-level-regulation}{%
\subsection{RQ3: Arc-Level Regulation}\label{rq3-arc-level-regulation}}

\textbf{Can regulatory logic (thresholds, Hill equations, inhibition) be
encoded directly on arcs?}

\begin{itemize}
\tightlist
\item
  \textbf{Classical Petri nets}: Arcs have weights (stoichiometry) but
  no regulatory semantics
\item
  \textbf{Biological regulation}: ATP inhibits phosphofructokinase when
  M(ATP) ≥ 5.0 mM (allosteric feedback)
\item
  \textbf{Current approach}: Implement inhibition in transition code
  (hidden from topology)
\item
  \textbf{Question}: Can arcs carry \textbf{threshold formulas} and
  \textbf{Hill equations} making regulation visible in network
  structure?
\end{itemize}

\textbf{Hypothesis}: Extending arcs with \textbf{type}
(normal/test/inhibitor) and \textbf{threshold functions} embeds
regulatory logic in topology, enabling visual reasoning and formal
analysis.

\hypertarget{rq4-atomic-conservation}{%
\subsection{RQ4: Atomic Conservation}\label{rq4-atomic-conservation}}

\textbf{Can biochemical formulas replace abstract tokens to enable
elemental balance analysis?}

\begin{itemize}
\tightlist
\item
  \textbf{Classical Petri nets}: Places hold tokens (abstract, countable
  units)
\item
  \textbf{Biochemistry}: Molecules have elemental composition (C₆H₁₂O₆ =
  glucose)
\item
  \textbf{Mass balance}: Reactions must conserve atoms (6 carbons in \textrightarrow 6
  carbons out)
\item
  \textbf{Question}: Can place names be biochemical formulas enabling
  automatic stoichiometry validation?
\end{itemize}

\textbf{Hypothesis}: Associating places with \textbf{formulas}
(element\textrightarrowcount dictionaries) enables \textbf{elemental balance
verification} at the Petri net level, detecting modeling errors
(unbalanced reactions).

\hypertarget{rq5-biological-validity}{%
\subsection{RQ5: Biological Validity}\label{rq5-biological-validity}}

\textbf{Does the extended formalism preserve biological correctness
principles?}

\begin{itemize}
\tightlist
\item
  \textbf{Enzyme conservation}: Catalysts are not consumed (test arcs
  must preserve marking)
\item
  \textbf{Superposition}: Multiple reactions produce the same metabolite
  (additive, not conflicting)
\item
  \textbf{Stoichiometric precision}: Integer coefficients, no fractional
  molecules in discrete models
\item
  \textbf{Regulatory consistency}: Inhibitor arcs block transitions
  consistently (threshold always enforced)
\item
  \textbf{Question}: Can we prove the formalism respects these
  biological constraints?
\end{itemize}

\textbf{Hypothesis}: \textbf{Well-formedness constraints} (formal rules
on network structure) enforce biological validity mechanically, ruling
out invalid models.

\hypertarget{rq6-practical-applicability}{%
\subsection{RQ6: Practical
Applicability}\label{rq6-practical-applicability}}

\textbf{Can real biological systems be successfully modeled with the
extended formalism?}

\begin{itemize}
\tightlist
\item
  \textbf{Validation requirement}: Theoretical formalism is only useful
  if practitioners can model actual systems
\item
  \textbf{Examples}: Glycolysis (10 reactions, 3 regulatory
  checkpoints), TCA cycle (8 reactions, cyclic), cellular respiration
  (32 reactions, spanning glycolysis + TCA + oxidative phosphorylation)
\item
  \textbf{Question}: Do the four innovations (weak independence,
  heterogeneous types, arc regulation, formulas) enable modeling complex
  real pathways?
\end{itemize}

\textbf{Hypothesis}: A \textbf{progressive example series} (16 models
from simple ATP hydrolysis to complete cellular respiration)
demonstrates practical feasibility and identifies formalism limitations.

\begin{center}\rule{0.5\linewidth}{0.5pt}\end{center}

\hypertarget{thesis-contributions}{%
\section{1.4 Thesis Contributions}\label{thesis-contributions}}

This thesis makes \textbf{theoretical, validation, and implementation
contributions} to systems biology modeling.

\hypertarget{theoretical-contributions-part-ii-core-theory}{%
\subsection{1.4.1 Theoretical Contributions (Part II: Core
Theory)}\label{theoretical-contributions-part-ii-core-theory}}

\textbf{Chapter 4: Extended Biological Petri Net Formalism} -
\textbf{12-tuple definition}: (P, T, F, W, M₀, K, Φ, Σ, Θ, Δ, τ, ρ) -
Classical components (P, T, F, W, M₀, K): Places, transitions, arcs,
weights, initial marking, capacities - \textbf{Extensions}: -
\textbf{Φ}: Rate functions (Michaelis-Menten, Hill, mass action, custom)
- \textbf{Σ}: Test arcs (non-consumptive, for catalysis) - \textbf{Θ}:
Inhibitor arcs (threshold-based blocking) - \textbf{Δ}: Threshold
functions (constant, dynamic, Hill equation) - \textbf{τ}: Transition
types (Continuous, Stochastic, Timed, Burst) - \textbf{ρ}: Biochemical
formulas (element\textrightarrowcount dictionaries) - \textbf{Arc semantics}: Firing
rules for normal/test/inhibitor arcs across all transition types -
\textbf{Well-formedness constraints}: 8 formal rules enforcing
biological validity (enzyme conservation, stoichiometric consistency)

\textbf{Chapter 5: Weak Independence Theory} - \textbf{Formal
definition}: Transitions t₁, t₂ are \textbf{weakly independent} if (•t₁
∩ •t₂) = ∅ (disjoint input places) - Allows shared outputs (convergent
reactions \textrightarrow superposition) - Allows shared catalysts via test arcs
(enzyme serves multiple reactions) - \textbf{Dependency classification
algorithm}: O(\textbar T\textbar² · \textbar P\textbar) algorithm
assigning CONFLICT, COUPLING, or INDEPENDENT to all transition pairs -
\textbf{Reachability preservation theorem}: Weakly independent
transitions can fire concurrently without affecting reachability set -
\textbf{Biological cooperativity}: Formalization of how metabolic
pathways cooperate (shared intermediates, shared enzymes)

\textbf{Chapter 6: Biochemical Formula Tracking} - \textbf{Formula
representation}: Hill notation (C₆H₁₂O₆) \textrightarrow element-count dictionary
\{``C'': 6, ``H'': 12, ``O'': 6\} - \textbf{Elemental balance
verification}: Algorithm checking ∑(atoms\_in) = ∑(atoms\_out) for all
transitions - \textbf{Elemental balance matrix} S\_e: (elements ×
transitions), each entry = net atom change - \textbf{Cofactor
suggestion}: Algorithm proposing missing cofactors (H₂O, H⁺, Pi) when
reactions are unbalanced - \textbf{Database integration}: KEGG/ChEBI
formula retrieval for automatic model enrichment

\hypertarget{validation-contributions-part-iii-empirical-validation}{%
\subsection{1.4.2 Validation Contributions (Part III: Empirical
Validation)}\label{validation-contributions-part-iii-empirical-validation}}

\textbf{Chapter 7: Validation Through Examples} - \textbf{16-example
progressive series}: From simple (ATP hydrolysis) to complex (complete
cellular respiration) - \textbf{Phase structure}: - Phase 1 (Examples
01-03): Foundation (basic reactions, reversibility, catalysis) - Phase 2
(Examples 04-06): Regulation (inhibitor arcs, competitive inhibition,
feedback) - Phase 3 (Examples 07-08): Integration (multi-step pathways,
heterogeneous types) - Phase 4 (Examples 09-13): Complete pathways
(glycolysis, TCA, respiration) - Phase 5 (Examples 14-16): Advanced
(branching, resource competition, dynamic thresholds) - \textbf{Key
proof}: Example 08 (Energy Sensing Motif) demonstrates \textbf{all four
innovations simultaneously}: - Weak independence: PFK and PK have
disjoint inputs - Heterogeneous types: Continuous (PFK, PK, ATPase) +
Stochastic burst (Gene\_PFK) - Arc regulation: ATP inhibitor arcs on PFK
and PK (thresholds 2.5 mM, 2.0 mM) - Atomic conservation: All reactions
elementally balanced (verified) - \textbf{Quantitative validation}:
Parameters from BRENDA database, simulation results match literature -
\textbf{Scalability analysis}: Parallel execution achieves 2-4× speedup
on typical networks (8 cores)

\hypertarget{implementation-contributions-part-iv-supporting-tools}{%
\subsection{1.4.3 Implementation Contributions (Part IV: Supporting
Tools)}\label{implementation-contributions-part-iv-supporting-tools}}

\textbf{Chapter 8: SHYpn System Architecture} - Three-tier architecture
(Presentation/Business Logic/Data layers) - Model representation:
Faithful 12-tuple implementation in Python - UI components: Network
canvas (Cairo rendering), property panels (GTK4) - Persistence: Native
JSON format + SBML import + GraphML export

\textbf{Chapter 9: KEGG Compound Integration} - REST API wrapper for
KEGG database (18,000+ compounds, 11,000+ reactions) - Automatic formula
retrieval: User enters KEGG ID \textrightarrow formula auto-filled - Reaction import:
KEGG reaction ID \textrightarrow creates places, transition, arcs automatically -
Pathway bulk import: Entire glycolysis pathway imported in \textless5
seconds - Cofactor suggestion: Algorithm proposes missing H₂O, H⁺, Pi
based on elemental imbalance - Local caching: SQLite database \textrightarrow 60×
speedup on repeated queries

\textbf{Chapter 10: Intelligent Parameter Inference from BRENDA} - SOAP
API integration with BRENDA (2.7 million kinetic parameters) -
Statistical aggregation: Median Km + 95\% confidence interval (robust to
outliers) - Context-aware heuristics: - Organism priority
(\emph{Saccharomyces cerevisiae} \textgreater{} human \textgreater{}
all) - Substrate fuzzy matching (``glucose'' matches ``D-glucose'',
``α-D-glucose'') - Quality filtering (citation presence, experimental
conditions, outlier detection) - Local database: Progressive
accumulation \textrightarrow 200× speedup on cached data - Batch enrichment: Entire
glycolysis pathway (10 enzymes) parameterized in \textless10 seconds

\textbf{Chapter 11: Hybrid Simulation Engine} - Four-engine
architecture: - Continuous: SciPy RK45 ODE solver (adaptive
time-stepping) - Stochastic: Gillespie SSA (exact algorithm) - Timed:
Priority queue for scheduled events - Burst: Exponential burst frequency
+ geometric burst size - Hybrid scheduler: Event-driven coordination
(asks all engines ``when is your next event?'') - Parallel execution:
Weak independence-based partitioning \textrightarrow 2-4× speedup (8 cores) -
Validation: Correctness proofs (reachability preservation), benchmark
suite (16 examples)

\hypertarget{impact-summary}{%
\subsection{1.4.4 Impact Summary}\label{impact-summary}}

\textbf{This thesis enables, for the first time}: 1. \textbf{Parallel
simulation exploiting biological cooperativity}: Weak independence
theory allows 2-4× speedup on multi-core systems by executing
non-conflicting reactions simultaneously 2. \textbf{Multi-scale
integrated models}: Single model spans continuous enzyme kinetics
(glycolysis) and stochastic gene expression bursts (\emph{lacZ}
transcription) 3. \textbf{Topology-embedded regulation}: Inhibitor arcs
with threshold formulas make regulatory logic visible in network
structure (no hidden code) 4. \textbf{Atomic-level mass balance}:
Biochemical formula tracking enables elemental conservation
verification, detecting modeling errors automatically

\textbf{Broader significance}: - \textbf{Computational systems biology}:
Provides formal foundation for integrated modeling (metabolomics +
transcriptomics) - \textbf{Synthetic biology}: Enables design and
simulation of engineered pathways with regulatory feedback -
\textbf{Drug discovery}: Models metabolic perturbations (enzyme
inhibition) coupled to gene expression responses - \textbf{Education}:
Visual formalism makes biochemical regulation accessible to students
(arc types visible)

\begin{center}\rule{0.5\linewidth}{0.5pt}\end{center}

\hypertarget{research-methodology}{%
\section{1.5 Research Methodology}\label{research-methodology}}

\hypertarget{formalism-development}{%
\subsection{1.5.1 Formalism Development}\label{formalism-development}}

\textbf{Iterative refinement approach}: 1. \textbf{Requirements
analysis}: Survey biological systems needing integrated modeling (lac
operon, glycolysis regulation, etc.) 2. \textbf{Formalism design}:
Extend Petri net 5-tuple to 12-tuple, adding Φ, Σ, Θ, Δ, τ, ρ components
3. \textbf{Formal specification}: Define enabling conditions, firing
rules, semantics for each transition type 4. \textbf{Theoretical
analysis}: Prove weak independence reachability preservation theorem 5.
\textbf{Validation}: Test formalism on 16 progressive examples (simple \textrightarrow
complex)

\textbf{Design criteria}: - \textbf{Minimality}: Add only necessary
components (no feature bloat) - \textbf{Compositionality}: Models
combine without side effects - \textbf{Biological fidelity}: Respects
fundamental biological principles (enzyme conservation, mass balance) -
\textbf{Formal analyzability}: Enables structural analysis
(P-invariants, reachability, weak independence classification)

\hypertarget{validation-strategy}{%
\subsection{1.5.2 Validation Strategy}\label{validation-strategy}}

\textbf{Progressive example series} (16 models): - \textbf{Phase 1}:
Establish baseline (simple reactions, reversibility, catalysis) -
\textbf{Phase 2}: Add regulation (inhibitor arcs, feedback loops) -
\textbf{Phase 3}: Demonstrate integration (continuous + stochastic, weak
independence) - \textbf{Phase 4}: Scale complexity (complete pathways:
glycolysis, TCA, respiration) - \textbf{Phase 5}: Stress-test limits
(branching, competition, dynamic thresholds)

\textbf{Validation criteria per example}: 1. \textbf{Structural
correctness}: Satisfies well-formedness constraints (C1-C8) 2.
\textbf{Elemental balance}: All transitions conserve atoms (verified
automatically) 3. \textbf{Parameter realism}: Km, Vmax from BRENDA
database (not arbitrary) 4. \textbf{Simulation convergence}: Reaches
steady-state or expected dynamics 5. \textbf{Literature consistency}:
Results match published experimental data

\textbf{Key proof example}: Example 08 (Energy Sensing Motif) is the
\textbf{minimal model demonstrating all four innovations}. If this
example works correctly (which it does, verified), it proves the
formalism is viable.

\hypertarget{implementation-validation}{%
\subsection{1.5.3 Implementation
Validation}\label{implementation-validation}}

\textbf{SHYpn platform serves as proof-of-concept}: - Implements all 12
formalism components - Executes all 16 validation examples successfully
- Achieves 2-4× speedup via parallel execution (8 cores) - Integrates
KEGG (formulas) and BRENDA (parameters)

\textbf{Implementation is secondary contribution}: Thesis proves
formalism is sound; tool proves formalism is executable.

\begin{center}\rule{0.5\linewidth}{0.5pt}\end{center}

\hypertarget{thesis-scope-and-limitations}{%
\section{1.6 Thesis Scope and
Limitations}\label{thesis-scope-and-limitations}}

\hypertarget{in-scope}{%
\subsection{1.6.1 In Scope}\label{in-scope}}

\textbf{What this thesis addresses}: - \textbf{Formalism}: Extending
Petri nets for integrated biological modeling - \textbf{Theory}: Weak
independence, heterogeneous transitions, arc regulation, formula
tracking - \textbf{Validation}: 16 biochemical examples demonstrating
practical applicability - \textbf{Implementation}: SHYpn platform as
proof-of-concept

\textbf{Biological domains covered}: - Central metabolism (glycolysis,
TCA cycle, oxidative phosphorylation) - Gene regulation (transcription,
mRNA dynamics, protein production) - Simple feedback loops (product
inhibition, feed-forward loops)

\hypertarget{out-of-scope}{%
\subsection{1.6.2 Out of Scope}\label{out-of-scope}}

\textbf{What this thesis does not address}: - \textbf{Spatial dynamics}:
Petri nets model well-mixed compartments (no diffusion, cell geometry) -
\textbf{Limitation}: Cannot model spatial gradients (e.g., morphogen
gradients in development) - \textbf{Future work}: Extend to spatially
distributed Petri nets

\begin{itemize}
\tightlist
\item
  \textbf{Multi-compartment systems}: Current formalism is
  single-compartment

  \begin{itemize}
  \tightlist
  \item
    \textbf{Limitation}: Cannot model cytoplasm-mitochondria
    interactions with transport
  \item
    \textbf{Future work}: Add compartment types, transport arcs
  \end{itemize}
\item
  \textbf{Protein-protein interactions}: Focus is metabolites + genes

  \begin{itemize}
  \tightlist
  \item
    \textbf{Limitation}: Signaling cascades (MAPK, JAK-STAT) not fully
    explored
  \item
    \textbf{Future work}: Extend to protein complexes,
    post-translational modifications
  \end{itemize}
\item
  \textbf{Evolutionary dynamics}: No mutation, selection, population
  genetics

  \begin{itemize}
  \tightlist
  \item
    \textbf{Out of scope}: Thesis is mechanistic modeling, not
    evolutionary
  \end{itemize}
\item
  \textbf{Whole-cell models}: Validated on pathways (10-50 transitions),
  not entire genomes

  \begin{itemize}
  \tightlist
  \item
    \textbf{Limitation}: Scalability to 1000+ transition networks
    unproven
  \item
    \textbf{Future work}: Hierarchical modeling, abstraction techniques
  \end{itemize}
\end{itemize}

\hypertarget{assumptions}{%
\subsection{1.6.3 Assumptions}\label{assumptions}}

\textbf{Modeling assumptions}: 1. \textbf{Well-mixed compartments}: No
spatial heterogeneity (stirred reactor assumption) 2.
\textbf{Deterministic enzyme concentrations}: {[}E{]}₀ is fixed (no
enzyme degradation modeled) 3. \textbf{Steady-state approximation for
complexes}: Enzyme-substrate complex formation is fast 4. \textbf{Ideal
solution}: No crowding effects, activity coefficients = 1 5.
\textbf{Constant pH/temperature}: Environmental parameters fixed

\textbf{These assumptions are standard in systems biology} and do not
invalidate the formalism's contributions.

\begin{center}\rule{0.5\linewidth}{0.5pt}\end{center}

\hypertarget{thesis-organization}{%
\section{1.7 Thesis Organization}\label{thesis-organization}}

\textbf{Part I: Introduction and Foundations} (Chapters 1-3) -
\textbf{Chapter 1}: Motivation, research questions, contributions (this
chapter) - \textbf{Chapter 2}: Background (Petri nets, biological
modeling, related work) - \textbf{Chapter 3}: Integration challenge (lac
operon deep dive, requirements analysis)

\textbf{Part II: Core Theory - The Extended Formalism} (Chapters 4-6) -
\textbf{Chapter 4}: Extended Bio-PN definition (12-tuple, arc semantics,
well-formedness) - \textbf{Chapter 5}: Weak independence theory
(definition, algorithm, reachability theorem) - \textbf{Chapter 6}:
Biochemical formula tracking (elemental balance, cofactor suggestion)

\textbf{Part III: Validation Through Examples} (Chapter 7) -
\textbf{Chapter 7}: 16-example progressive series (Phase 1-5,
quantitative validation)

\textbf{Part IV: Implementation - Supporting Tools} (Chapters 8-11) -
\textbf{Chapter 8}: SHYpn system architecture (three-tier design, model
representation) - \textbf{Chapter 9}: KEGG integration (formula
retrieval, reaction import, caching) - \textbf{Chapter 10}: Parameter
inference (BRENDA API, statistical aggregation, heuristics) -
\textbf{Chapter 11}: Hybrid simulation engine (four engines, scheduler,
parallel execution)

\textbf{Part V: Evaluation} (Chapters 12-13) - \textbf{Chapter 12}: Case
studies (glycolysis regulation, TCA cycle, cellular respiration) -
\textbf{Chapter 13}: Performance evaluation (scalability, speedup
analysis, comparison)

\textbf{Part VI: Synthesis} (Chapters 14-15) - \textbf{Chapter 14}:
Discussion (contributions, limitations, biological validity) -
\textbf{Chapter 15}: Conclusion and future work (research impact, open
problems)

\textbf{Estimated length}: 180-200 pages (10-12 pages per chapter
average)

\begin{center}\rule{0.5\linewidth}{0.5pt}\end{center}

\hypertarget{reading-guide}{%
\section{1.8 Reading Guide}\label{reading-guide}}

\textbf{For readers interested in}:

\textbf{Theoretical foundations}: - Read Chapters 4-6 (formalism
definition, weak independence, formula tracking) - Skip implementation
details (Chapters 8-11)

\textbf{Practical modeling}: - Read Chapter 7 (validation examples) -
Read Chapters 9-10 (KEGG/BRENDA integration for parameter inference) -
Skim theory (Chapters 4-6) for terminology

\textbf{Implementation and tools}: - Read Chapters 8-11 (architecture,
simulation engine, parallel execution) - Chapter 7 for example suite

\textbf{Systems biology applications}: - Read Chapter 3 (integration
challenge, lac operon) - Read Chapter 7 (validation examples) - Read
Chapter 12 (case studies)

\textbf{Computer science formalism}: - Read Chapters 4-5 (12-tuple
definition, weak independence algorithm) - Chapter 11 (hybrid simulation
semantics)

\begin{center}\rule{0.5\linewidth}{0.5pt}\end{center}

\hypertarget{summary}{%
\section{1.9 Summary}\label{summary}}

\textbf{This chapter established the need for integrated biological
modeling}, motivated by the lac operon system spanning metabolism and
gene regulation. \textbf{Six research questions} guide the thesis: 1.
Can Petri nets support parallel execution with shared places?
(\textbf{RQ1: Weak independence}) 2. Can continuous, stochastic, timed,
and burst transitions coexist? (\textbf{RQ2: Heterogeneity}) 3. Can
regulatory logic be encoded on arcs? (\textbf{RQ3: Arc regulation}) 4.
Can biochemical formulas enable atomic conservation? (\textbf{RQ4:
Formula tracking}) 5. Does the formalism preserve biological
correctness? (\textbf{RQ5: Validity}) 6. Can real systems be modeled?
(\textbf{RQ6: Applicability})

\textbf{Contributions span theory} (Extended Bio-PN formalism, weak
independence theory), \textbf{validation} (16-example progressive
series), and \textbf{implementation} (SHYpn platform with KEGG/BRENDA
integration).

\textbf{Chapter 2 surveys related work}, establishing context for this
thesis's innovations. \textbf{Chapter 3 presents the integration
challenge} in depth, deriving formal requirements for unified modeling.
